% Options for packages loaded elsewhere
\PassOptionsToPackage{unicode}{hyperref}
\PassOptionsToPackage{hyphens}{url}
%
\documentclass[
]{article}
\usepackage{amsmath,amssymb}
\usepackage{iftex}
\ifPDFTeX
  \usepackage[T1]{fontenc}
  \usepackage[utf8]{inputenc}
  \usepackage{textcomp} % provide euro and other symbols
\else % if luatex or xetex
  \usepackage{unicode-math} % this also loads fontspec
  \defaultfontfeatures{Scale=MatchLowercase}
  \defaultfontfeatures[\rmfamily]{Ligatures=TeX,Scale=1}
\fi
\usepackage{lmodern}
\ifPDFTeX\else
  % xetex/luatex font selection
\fi
% Use upquote if available, for straight quotes in verbatim environments
\IfFileExists{upquote.sty}{\usepackage{upquote}}{}
\IfFileExists{microtype.sty}{% use microtype if available
  \usepackage[]{microtype}
  \UseMicrotypeSet[protrusion]{basicmath} % disable protrusion for tt fonts
}{}
\makeatletter
\@ifundefined{KOMAClassName}{% if non-KOMA class
  \IfFileExists{parskip.sty}{%
    \usepackage{parskip}
  }{% else
    \setlength{\parindent}{0pt}
    \setlength{\parskip}{6pt plus 2pt minus 1pt}}
}{% if KOMA class
  \KOMAoptions{parskip=half}}
\makeatother
\usepackage{xcolor}
\usepackage{longtable,booktabs,array}
\usepackage{calc} % for calculating minipage widths
% Correct order of tables after \paragraph or \subparagraph
\usepackage{etoolbox}
\makeatletter
\patchcmd\longtable{\par}{\if@noskipsec\mbox{}\fi\par}{}{}
\makeatother
% Allow footnotes in longtable head/foot
\IfFileExists{footnotehyper.sty}{\usepackage{footnotehyper}}{\usepackage{footnote}}
\makesavenoteenv{longtable}
\setlength{\emergencystretch}{3em} % prevent overfull lines
\providecommand{\tightlist}{%
  \setlength{\itemsep}{0pt}\setlength{\parskip}{0pt}}
\setcounter{secnumdepth}{-\maxdimen} % remove section numbering
\ifLuaTeX
  \usepackage{selnolig}  % disable illegal ligatures
\fi
\IfFileExists{bookmark.sty}{\usepackage{bookmark}}{\usepackage{hyperref}}
\IfFileExists{xurl.sty}{\usepackage{xurl}}{} % add URL line breaks if available
\urlstyle{same}
\hypersetup{
  hidelinks,
  pdfcreator={LaTeX via pandoc}}

\author{}
\date{}

\begin{document}

TRƯỜNG ĐẠI HỌC BÁCH KHOA HÀ NỘI

VIỆN CÔNG NGHỆ THÔNG TIN VÀ TRUYỀN THÔNG

──────── * ───────

\textbf{BÁO CÁO}

MÔN: Project 2

\textbf{Nghiên cứu và xây dựng trình dịch ngược mã máy thành lệnh hợp
ngữ của kiến trúc RISC-V}

\begin{quote}
Mã lớp học : 761986

Giáo viên hướng dẫn : Ngô Lam Trung

Sinh viên thực hiện:
\end{quote}

\begin{longtable}[]{@{}
  >{\raggedright\arraybackslash}p{(\columnwidth - 2\tabcolsep) * \real{0.2950}}
  >{\raggedright\arraybackslash}p{(\columnwidth - 2\tabcolsep) * \real{0.7050}}@{}}
\toprule\noalign{}
\begin{minipage}[b]{\linewidth}\raggedright
\textbf{Họ và Tên}
\end{minipage} & \begin{minipage}[b]{\linewidth}\raggedright
: Đinh Đức Mạnh
\end{minipage} \\
\midrule\noalign{}
\endhead
\bottomrule\noalign{}
\endlastfoot
\textbf{MSSV} & : 20235371 \\
\textbf{Email} & : Manh.DD235371@sis.hust.edu.vn \\
\end{longtable}

\emph{\textbf{Hà Nội, tháng 3 năm 2026}}

\textbf{MỤC LỤC}

\textbf{LỜI NÓI ĐẦU} 3

\textbf{CHƯƠNG 1. MỞ ĐẦU} 4

\textbf{CHƯƠNG 2. CƠ SỞ LÝ THUYẾT VÀ CÔNG NGHỆ SỬ DỤNG} 6

\textbf{CHƯƠNG 3. PHÂN TÍCH THIẾT KẾ HỆ THỐNG} 11

\textbf{CHƯƠNG 4. HƯỚNG DẪN TRIỂN KHAI VÀ CÀI ĐẶT} 12

\textbf{CHƯƠNG 5. KẾT LUẬN} 13

\textbf{TÀI LIỆU THAM KHẢO} 14

\textbf{LỜI NÓI ĐẦU}

Trong bối cảnh các hệ thống phần mềm ngày càng phức tạp và phụ thuộc sâu
vào kiến trúc phần cứng, việc nghiên cứu chương trình ở mức mã máy trở
thành nhu cầu quan trọng đối với kiểm thử, gỡ lỗi, phân tích bảo mật và
đánh giá hành vi thực thi. Nhiều hành vi của chương trình chỉ thực sự
được bộc lộ đầy đủ khi quan sát trực tiếp ở tầng instruction.

Đối với các tệp thực thi không có mã nguồn, các mẫu firmware, object
file hoặc các trường hợp phục vụ phân tích mã độc, nhà nghiên cứu không
thể đọc trực tiếp dãy byte nhị phân để hiểu ngữ nghĩa chương trình. Khi
đó, disassembler đóng vai trò như công cụ trung gian, giúp chuyển dữ
liệu nhị phân thành biểu diễn hợp ngữ gần với cách con người và công cụ
phân tích có thể tiếp cận.

Xuất phát từ nhu cầu trên, báo cáo này tập trung nghiên cứu và xây dựng
một trình dịch ngược cho kiến trúc RISC-V, cụ thể là biến thể RV32I. Mục
tiêu của đề tài không chỉ là hiển thị mã hợp ngữ, mà còn là thiết kế một
lõi xử lý có cấu trúc dữ liệu rõ ràng, có khả năng làm nền cho các bước
phát triển tiếp theo như phân tích luồng điều khiển, phục hồi nhãn, và
tích hợp với các công cụ dịch ngược bậc cao hơn.

Trong phạm vi hiện tại, báo cáo hoàn thiện phần mở đầu, cơ sở lý thuyết,
phân tích công cụ tham khảo và định hướng kiến trúc phần mềm. Các chương
sau về phân tích thiết kế chi tiết, triển khai và kết luận được giữ dưới
dạng placeholder để tiếp tục phát triển ở các giai đoạn sau.

\hypertarget{mux1edf-ux111ux1ea7u}{%
\section{MỞ ĐẦU}\label{mux1edf-ux111ux1ea7u}}

\hypertarget{luxfd-do-chux1ecdn-ux111ux1ec1-tuxe0i}{%
\subsection{\texorpdfstring{ Lý do chọn đề
tài}{ Lý do chọn đề tài}}\label{luxfd-do-chux1ecdn-ux111ux1ec1-tuxe0i}}

Máy tính hiện đại được tổ chức thành nhiều tầng trừu tượng: phần cứng,
hệ điều hành, trình biên dịch, thư viện, chương trình và giao diện người
dùng. Tuy nhiên, giữa ngôn ngữ nguồn và tầng thực thi luôn tồn tại một
khoảng cách nhận thức đáng kể. Nhiều lỗi hệ thống, lỗ hổng bảo mật hoặc
hành vi bất thường chỉ bộc lộ rõ khi quan sát chương trình ở mức mã máy.

Từ góc độ an ninh và phân tích nhị phân, disassembler là công cụ nền
tảng trong ba bối cảnh tiêu biểu: phân tích mã độc, nghiên cứu lỗ hổng
bảo mật và pháp y kỹ thuật số. Trong các tình huống này, đối tượng phân
tích thường chỉ tồn tại dưới dạng file thực thi hoặc object file; vì vậy
cần một công cụ có khả năng đọc, giải mã và tái hiện dòng lệnh hợp ngữ
để phục vụ việc lập luận kỹ thuật.

\begin{quote}
• Phân tích mã độc: hỗ trợ bóc tách các mẫu chương trình độc hại không
có mã nguồn, từ đó xác định hành vi, luồng điều khiển và dấu hiệu né
tránh phân tích.

• Nghiên cứu lỗ hổng bảo mật: cho phép kiểm tra các đoạn mã nhị phân của
thư viện, firmware hoặc ứng dụng đóng để phát hiện logic xử lý nguy
hiểm.

• Pháp y kỹ thuật số: hỗ trợ trích xuất bằng chứng từ phần mềm thu giữ
được trong quá trình điều tra, nhất là khi tác giả phần mềm không cung
cấp thông tin triển khai.
\end{quote}

Mặc dù GNU objdump và nhiều công cụ mature khác đã hỗ trợ disassembly
cho RISC-V, chúng thường được thiết kế theo hướng đa kiến trúc, đa định
dạng và ưu tiên đầu ra văn bản cho người đọc. Đồ án này lựa chọn cách
tiếp cận tối giản hơn: xây dựng một lõi chuyên biệt cho RV32I, tập trung
vào tính đúng đắn của parser, resolver, decoder và mô hình dữ liệu đầu
ra để phục vụ học tập cũng như nghiên cứu sâu hơn.

\hypertarget{mux1ee5c-tiuxeau-ux111ux1ec1-tuxe0i}{%
\subsection{\texorpdfstring{ Mục tiêu đề
tài}{ Mục tiêu đề tài}}\label{mux1ee5c-tiuxeau-ux111ux1ec1-tuxe0i}}

\hypertarget{mux1ee5c-tiuxeau-tux1ed5ng-quuxe1t}{%
\subsubsection{\texorpdfstring{ Mục tiêu tổng
quát}{ Mục tiêu tổng quát}}\label{mux1ee5c-tiuxeau-tux1ed5ng-quuxe1t}}

Nghiên cứu nguyên lý hoạt động của disassembler và xây dựng một công cụ
có khả năng đọc tệp nhị phân RISC-V, xác định vùng chứa mã lệnh, giải mã
các instruction RV32I và xuất ra biểu diễn hợp ngữ tương ứng với độ
trung thực cao đối với chương trình gốc.

\hypertarget{mux1ee5c-tiuxeau-cux1ee5-thux1ec3}{%
\subsubsection{\texorpdfstring{ Mục tiêu cụ
thể}{ Mục tiêu cụ thể}}\label{mux1ee5c-tiuxeau-cux1ee5-thux1ec3}}

\begin{quote}
• Tìm hiểu cơ sở lý thuyết của disassembler và vai trò của disassembly
trong phân tích chương trình ở mức thấp.

• Nghiên cứu kiến trúc tập lệnh RISC-V, đặc biệt là RV32I, bao gồm định
dạng lệnh, thanh ghi, opcode, funct3, funct7 và cơ chế mã hóa immediate.

• Nghiên cứu cấu trúc file ELF và object file liên quan tới quá trình
đọc section, symbol và relocation cần thiết cho disassembly.

• Xây dựng parser để đọc ELF đầu vào, xác định section mã và ánh xạ địa
chỉ logic của instruction.

• Xây dựng bộ giải mã RV32I để biến dãy byte thành instruction có cấu
trúc và có thể phát sinh đầu ra ở nhiều định dạng.

• Thiết kế structured output làm nền cho các bước mở rộng như dựng CFG,
tích hợp decompiler hoặc kiểm thử đối chiếu với công cụ chuẩn.

• Đối chiếu kết quả với công cụ tham khảo như GNU objdump để đánh giá độ
đúng của hệ thống.
\end{quote}

\hypertarget{phux1ea1m-vi-ux111ux1ec1-tuxe0i}{%
\subsubsection{\texorpdfstring{ Phạm vi đề
tài}{ Phạm vi đề tài}}\label{phux1ea1m-vi-ux111ux1ec1-tuxe0i}}

Để đảm bảo tính khả thi trong khuôn khổ học phần, phạm vi triển khai
hiện tại của đề tài được giới hạn như sau:

\begin{quote}
• Đầu vào là các file nhị phân RISC-V phổ biến trong bài toán dịch ngược
cơ bản, ưu tiên ELF executable và relocatable object file (.o).

• Kiến trúc hỗ trợ chính là RISC-V RV32I; các extension như M, C, F, D
hoặc V chưa phải trọng tâm của giai đoạn hiện tại.

• Giao diện sử dụng là CLI để giảm độ phức tạp triển khai; giao diện đồ
họa và các tiện ích nâng cao được xem là hướng mở rộng.

• Phần lõi tập trung vào parser, resolver và decoder; các bài toán nâng
cao như function recovery, decompiler IR hoặc symbolic execution chưa
được triển khai đầy đủ trong báo cáo này.

• Đầu ra chuẩn được định hướng ở dạng JSON/structured output, đồng thời
có thể bổ sung lớp text emitter để đối chiếu với objdump khi cần.
\end{quote}

\hypertarget{cux1a1-sux1edf-luxfd-thuyux1ebft-vuxe0-cuxf4ng-nghux1ec7-sux1eed-dux1ee5ng}{%
\section{CƠ SỞ LÝ THUYẾT VÀ CÔNG NGHỆ SỬ
DỤNG}\label{cux1a1-sux1edf-luxfd-thuyux1ebft-vuxe0-cuxf4ng-nghux1ec7-sux1eed-dux1ee5ng}}

\hypertarget{cux1a1-sux1edf-luxfd-thuyux1ebft}{%
\subsection{\texorpdfstring{ Cơ sở lý
thuyết}{ Cơ sở lý thuyết}}\label{cux1a1-sux1edf-luxfd-thuyux1ebft}}

\hypertarget{tux1ed5ng-quan-vux1ec1-kiux1ebfn-truxfac-risc-v-vuxe0-biux1ebfn-thux1ec3-rv32i}{%
\subsubsection{\texorpdfstring{ Tổng quan về kiến trúc RISC-V và biến
thể
RV32I}{ Tổng quan về kiến trúc RISC-V và biến thể RV32I}}\label{tux1ed5ng-quan-vux1ec1-kiux1ebfn-truxfac-risc-v-vuxe0-biux1ebfn-thux1ec3-rv32i}}

RISC-V là kiến trúc tập lệnh mở theo triết lý RISC, được thiết kế để đơn
giản, mô-đun và thuận lợi cho cả nghiên cứu lẫn triển khai công nghiệp.
Trong hệ sinh thái RISC-V, RV32I là tập lệnh số nguyên cơ sở 32-bit, đủ
tối thiểu để làm đích biên dịch cho trình biên dịch và môi trường hệ
điều hành hiện đại {[}1{]}.

Ở mức lập trình, RV32I sử dụng 32 thanh ghi tổng quát x0 đến x31 cùng
thanh ghi đếm chương trình pc. Các instruction có độ dài cố định 32 bit
và được mã hóa theo các định dạng chính như R, I, S, B, U và J. Việc các
trường rd, rs1 và rs2 được đặt ở vị trí tương đối ổn định giữa các định
dạng giúp đơn giản hóa phần cứng giải mã, đồng thời hỗ trợ xây dựng
decoder phần mềm rõ ràng hơn {[}1{]}.

Đối với disassembler, các đặc tính này có hai ý nghĩa quan trọng. Thứ
nhất, việc mã hóa cố định 32 bit cho phép quét tuyến tính dễ dàng trong
phạm vi RV32I cơ sở. Thứ hai, các quy tắc sign-extension và bố trí bit
của immediate buộc decoder phải tách bit một cách chính xác thay vì chỉ
in theo khuôn mẫu chuỗi ký tự.

\hypertarget{cux1ea5u-truxfac-file-elf-vuxe0-uxfd-nghux129a-ux111ux1ed1i-vux1edbi-disassembler}{%
\subsubsection{\texorpdfstring{ Cấu trúc file ELF và ý nghĩa đối với
disassembler}{ Cấu trúc file ELF và ý nghĩa đối với disassembler}}\label{cux1ea5u-truxfac-file-elf-vuxe0-uxfd-nghux129a-ux111ux1ed1i-vux1edbi-disassembler}}

ELF (Executable and Linking Format) là định dạng đối tượng và thực thi
phổ biến trong các hệ thống UNIX-like. Theo đặc tả gABI, phần ELF header
đóng vai trò ``bản đồ'' mô tả cách tổ chức file; còn section giữ phần
lớn thông tin ở góc nhìn liên kết như mã lệnh, dữ liệu, bảng ký hiệu và
relocation {[}7{]}.

Trong bài toán disassembly, parser cần đọc chính xác section header
table để xác định vùng mã lệnh, nhất là .text, đồng thời tra cứu symbol
table để phục hồi tên hàm hoặc nhãn khi dữ liệu còn tồn tại. Nếu làm
việc với relocatable object file, thông tin relocation cũng quan trọng
vì nó cho biết những vị trí nào trong mã cần được diễn giải theo ngữ
cảnh liên kết chứ không chỉ là immediate thuần túy {[}8{]}.

Do đó, việc hiểu ELF không chỉ phục vụ bước ``đọc file'', mà còn quyết
định mức độ giàu ngữ nghĩa của output. Một disassembler chỉ quét byte
trong .text sẽ cho ra chuỗi assembly cơ bản; ngược lại, một disassembler
biết khai thác symbol và relocation có thể tạo ra đầu ra gần với ngữ
nghĩa chương trình hơn.

\hypertarget{phuxe2n-tuxedch-cuxf4ng-cux1ee5-tham-khux1ea3o-vuxe0-phux1b0ux1a1ng-phuxe1p-tiux1ebfp-cux1eadn-gnu-objdump}{%
\subsubsection{\texorpdfstring{ Phân tích công cụ tham khảo và phương
pháp tiếp cận: GNU
objdump}{ Phân tích công cụ tham khảo và phương pháp tiếp cận: GNU objdump}}\label{phuxe2n-tuxedch-cuxf4ng-cux1ee5-tham-khux1ea3o-vuxe0-phux1b0ux1a1ng-phuxe1p-tiux1ebfp-cux1eadn-gnu-objdump}}

GNU objdump là một trong những công cụ disassembly phổ biến nhất và là
case study phù hợp để rút ra phương pháp thiết kế. Về bản chất, objdump
thể hiện rõ tư tưởng phân lớp: lớp frontend nhận tham số dòng lệnh và
điều phối, BFD trừu tượng hóa định dạng tệp, còn libopcodes đảm nhiệm
logic giải mã theo từng kiến trúc {[}2{]}, {[}3{]}.

BFD sử dụng các cấu trúc canonical nội bộ để biểu diễn object file, nhờ
đó lớp giải mã không phải làm việc trực tiếp với từng định dạng nhị phân
cụ thể. Đây là ý tưởng có giá trị lớn đối với đề tài: nếu parser ELF và
decoder bị dính chặt vào nhau, hệ thống sẽ khó mở rộng và khó kiểm thử.
Ngược lại, một lớp chuẩn hóa nội bộ sẽ giúp cô lập phần phụ thuộc định
dạng khỏi lõi giải mã {[}3{]}.

Tuy nhiên, objdump cũng bộc lộ hai hạn chế khi nhìn từ mục tiêu của đồ
án này. Thứ nhất, đầu ra chính của nó là chuỗi văn bản nhằm phục vụ quan
sát của con người. Thứ hai, hệ sinh thái libbfd/libopcodes rất mạnh
nhưng khá đồ sộ nếu chỉ cần một công cụ chuyên biệt cho RV32I. Vì vậy,
báo cáo này chọn cách học theo tư tưởng kiến trúc của objdump nhưng giản
lược phạm vi hiện thực {[}2{]}--{[}4{]}.

\textbf{2.2. Định hướng kiến trúc phần mềm đề xuất}

\textbf{2.2.1. Nguyên tắc phân lớp và luồng xử lý}

Qua quá trình trao đổi và rà soát yêu cầu, kiến trúc phần mềm của hệ
thống được thống nhất theo hướng phân lớp rõ ràng. Từ góc nhìn triển
khai, hệ thống có thể được hiểu theo năm lớp trách nhiệm chính: lớp giao
diện CLI; lớp điều phối luồng xử lý; lớp truy cập dữ liệu nhị phân; lớp
lõi giải mã; và lớp mô hình dữ liệu cùng hợp đồng output.

Cụ thể hơn, presentation layer chịu trách nhiệm nhận input, option và
phản hồi cho người dùng. Application/orchestration layer điều phối
pipeline mở file, parse, resolve, decode và emit. Infrastructure layer
xử lý ELF loader/parser, byte access, endianness và các bảng header.
Core/domain layer chứa resolver, decoder RV32I và các logic quy tắc
nghiệp vụ. Cuối cùng, model/output layer biểu diễn BinaryImage,
SectionRecord, SymbolRecord, InstructionIR và các đối tượng dùng để
serialization.

Luồng xử lý tương ứng là: nhận file đầu vào → chuẩn hóa thông tin nhị
phân → xác định section mã → giải mã instruction → sinh structured
output → tùy chọn render thành văn bản assembly. Cách chia này bảo đảm
mỗi khối có một nhiệm vụ rõ ràng và giảm phụ thuộc chéo giữa parser và
decoder.

\textbf{2.2.2. Chuẩn hóa dữ liệu đầu vào và mô hình canonical nội bộ}

Một điểm đã được thống nhất là hệ thống nên học theo tư tưởng của BFD
nhưng không sao chép toàn bộ BFD. Nói cách khác, đề tài nên xây dựng một
mô hình canonical nội bộ kiểu ``mini-BFD'' thay vì cố hiện thực một
framework tổng quát cho nhiều định dạng và nhiều kiến trúc.\footnote{Trong
  báo cáo này, mô hình ``mini-BFD'' được hiểu là một lớp chuẩn hóa nội
  bộ tối giản cho ELF/.o, chỉ giữ những trường thật sự cần cho pipeline
  disassembly và không nhằm thay thế đầy đủ thư viện BFD của GNU
  Binutils.}

Trong mô hình này, ELF parser không trả trực tiếp các byte rời rạc cho
decoder, mà chuyển chúng về các cấu trúc chuẩn hóa như BinaryImage,
SectionRecord, SymbolRecord và RelocationRecord. Từ đó, decoder RV32I
chỉ cần làm việc với code section bytes, base address và ánh xạ symbol
cần thiết, không cần hiểu chi tiết của toàn bộ ELF header trong mọi bước
xử lý.

Cách tiếp cận này vừa giữ lại ưu điểm lớn nhất của BFD --- tách parser
khỏi logic giải mã --- vừa tránh được độ cồng kềnh không cần thiết đối
với một đồ án tập trung vào RV32I. Đây là điểm cân bằng hợp lý giữa giá
trị học thuật và tính khả thi triển khai.

\textbf{2.2.3. Instruction IR và structured output}

Đầu ra của hệ thống không nên dừng ở chuỗi assembly thuần túy. Thay vào
đó, mỗi instruction cần được lưu dưới dạng một biểu diễn trung gian có
cấu trúc, gọi là Instruction IR, bao gồm các trường tối thiểu như địa
chỉ, kích thước, bytes gốc, mnemonic, operands, thanh ghi đọc/ghi,
branch target và symbol liên quan.\footnote{Instruction IR (Intermediate
  Representation) là biểu diễn trung gian có cấu trúc của một lệnh máy
  sau khi giải mã; nó khác với chuỗi assembly ở chỗ được tối ưu cho phần
  mềm xử lý tiếp như CFG builder, data-flow analyzer hoặc decompiler.}

Việc đưa thêm lớp Instruction IR giúp tách biệt ``nghĩa của lệnh'' với
``cách hiển thị lệnh''. Từ cùng một nguồn dữ liệu, hệ thống có thể vừa
sinh chuỗi .asm cho người đọc, vừa sinh JSON cho công cụ xử lý tiếp, vừa
làm đầu vào cho các mô-đun phân tích luồng điều khiển hoặc phục hồi hàm.
Định hướng này cũng nhất quán với mục tiêu phát triển một lõi
disassembler có giá trị tái sử dụng chứ không chỉ là chương trình in văn
bản {[}9{]}.

Trong ngắn hạn, structured output được chọn là JSON vì JSON dễ sinh, dễ
kiểm tra bằng test và dễ chuyển đổi sang các định dạng khác. Trong trung
hạn, cùng một mô hình dữ liệu này có thể được serialize sang các format
chặt hơn như MessagePack hoặc Protocol Buffers nếu yêu cầu hiệu năng
hoặc tích hợp hệ thống tăng lên.

\begin{enumerate}
\def\labelenumi{\arabic{enumi}.}
\item
  \textbf{1. Định dạng file ELF}
\end{enumerate}

ELF (Executable and Linkable Format) là định dạng tệp chuẩn dùng để lưu
trữ chương trình thực thi, object file và thư viện. Trong bài toán
disassembler, việc hiểu đúng cấu trúc ELF là điều kiện tiên quyết để
chương trình có thể xác định chính xác vùng mã lệnh cần dịch. Một file
ELF thường gồm các thành phần quan trọng như \textbf{ELF Header},
\textbf{Program Header Table}, \textbf{Section Header Table} và các
\textbf{section} dữ liệu. Trong đó, ELF Header chứa các thông tin nhận
dạng file như magic number, kiến trúc mục tiêu, entry point và vị trí
các bảng mô tả; Section Header Table dùng để mô tả các section như
.text, .data, .bss, .symtab, .strtab; còn section .text thường là nơi
chứa mã máy cần disassemble. Vì vậy, muốn disassembler hoạt động đúng
thì trước hết phải đọc đúng các trường trong ELF, xác định đúng tên
section, offset, kích thước, địa chỉ và phạm vi dữ liệu của từng
section, đặc biệt là section chứa mã lệnh.

\begin{enumerate}
\def\labelenumi{\arabic{enumi}.}
\setcounter{enumi}{1}
\item
  \textbf{2. Cơ sở kiểm thử trình disassembler}
\end{enumerate}

Cơ sở kiểm thử của trình disassembler được xây dựng dựa trên việc đối
chiếu với chuẩn kỹ thuật và công cụ tham chiếu. Trước hết, chương trình
phải tuân thủ đúng \textbf{đặc tả ELF} khi đọc file và đúng \textbf{đặc
tả ISA} khi giải mã instruction. Tiếp theo, kết quả của chương trình cần
được so sánh với các công cụ chuẩn như readelf và objdump để kiểm tra
tính chính xác của việc đọc section, giải mã lệnh, xác định địa chỉ và
hiển thị toán hạng. Ngoài ra, cần xây dựng bộ test gồm các file ELF đơn
giản, các chương trình assembly nhỏ, các trường hợp biên như immediate
âm/dương, branch offset, jump target, cũng như các file lỗi hoặc file
ELF không hợp lệ để kiểm tra khả năng xử lý ngoại lệ. Từ đó có thể đánh
giá rằng chương trình đúng khi nó đọc đúng cấu trúc ELF, disassemble
đúng vùng mã, giải mã đúng instruction và cho kết quả nhất quán với công
cụ chuẩn trên các bộ dữ liệu kiểm thử.

\textbf{2.2.4. Đánh giá lựa chọn công nghệ triển khai}

Đối với phạm vi Project 2, Java 21 là lựa chọn hợp lý vì tốc độ phát
triển nhanh, mô hình đối tượng rõ ràng, thư viện chuẩn đủ mạnh để thao
tác file nhị phân và hệ sinh thái test/build trưởng thành. Maven hỗ trợ
quản lý dependency, chuẩn hóa lifecycle build và tăng tính lặp lại của
quá trình triển khai {[}5{]}, {[}6{]}.

So với C/C++, Java bất lợi ở chỗ khó tận dụng trực tiếp các thư viện
dịch ngược native vốn đã phổ biến trong cộng đồng reverse engineering.
Tuy nhiên, với mục tiêu hiện tại là nắm chắc parser, resolver, decoder
và mô hình dữ liệu, Java vẫn là lựa chọn cân bằng hơn. Hơn nữa, Java 21
đã cung cấp FFM API theo hướng gọi native library và thao tác native
memory thuận tiện hơn JNI truyền thống, vì vậy kiến trúc hiện tại hoàn
toàn có thể mở đường cho backend native ở giai đoạn sau {[}5{]}.

Từ góc nhìn thiết kế, quyết định hợp lý nhất là: Java được dùng để triển
khai lõi học thuật và structured output trong giai đoạn hiện tại; còn
khả năng tích hợp Capstone hoặc các thư viện native khác chỉ nên được
xem là hướng tối ưu hóa hoặc mở rộng về sau {[}10{]}.

\textbf{2.3. Công nghệ sử dụng}

Các công nghệ được lựa chọn ở giai đoạn hiện tại của hệ thống bao gồm:

\begin{quote}
• Java 21: ngôn ngữ lập trình chính để hiện thực parser, resolver,
decoder và mô hình dữ liệu.

• Apache Maven: công cụ quản lý dependency và chuẩn hóa quá trình
build/test/package của dự án.

• JUnit (dự kiến): bộ công cụ kiểm thử đơn vị để đối chiếu kết quả giải
mã với các trường hợp đầu vào chuẩn.

• GNU objdump: công cụ tham chiếu dùng cho việc kiểm thử đối chiếu kết
quả disassembly ở mức văn bản.

• JSON: định dạng output chuẩn ở giai đoạn đầu để phục vụ cả hiển thị
lẫn xử lý tiếp bằng công cụ khác.
\end{quote}

\textbf{2.4. Tóm tắt các nội dung đã thống nhất}

Bảng dưới đây tổng hợp ngắn gọn các quyết định kỹ thuật đã được thống
nhất trong quá trình trao đổi, đồng thời đóng vai trò như đầu vào cho
Chương 3 về phân tích thiết kế hệ thống.

\begin{longtable}[]{@{}
  >{\raggedright\arraybackslash}p{(\columnwidth - 4\tabcolsep) * \real{0.3095}}
  >{\raggedright\arraybackslash}p{(\columnwidth - 4\tabcolsep) * \real{0.3307}}
  >{\raggedright\arraybackslash}p{(\columnwidth - 4\tabcolsep) * \real{0.3598}}@{}}
\toprule\noalign{}
\begin{minipage}[b]{\linewidth}\raggedright
\textbf{Nội dung}
\end{minipage} & \begin{minipage}[b]{\linewidth}\raggedright
\textbf{Thống nhất}
\end{minipage} & \begin{minipage}[b]{\linewidth}\raggedright
\textbf{Ý nghĩa}
\end{minipage} \\
\midrule\noalign{}
\endhead
\bottomrule\noalign{}
\endlastfoot
Mô hình đầu vào & Chuẩn hóa ELF/.o sang mô hình canonical nội bộ tối
giản, không sao chép BFD đầy đủ. & Giảm phụ thuộc, tách parser khỏi
decoder và vẫn giữ khả năng mở rộng hợp lý. \\
Kiến trúc phần mềm & Phân tách giao diện, điều phối, truy cập nhị phân,
lõi giải mã và mô hình dữ liệu. & Giúp kiểm thử từng phần và thuận lợi
khi mở rộng sang CFG hoặc decompiler backend. \\
Ngôn ngữ & Java 21 + Maven cho giai đoạn hiện tại; để ngỏ khả năng tích
hợp native library về sau. & Cân bằng giữa tốc độ phát triển, độ an toàn
bộ nhớ và tính tổ chức của mã nguồn. \\
Đầu ra & JSON/structured output là đầu ra chuẩn; text emitter chỉ là lớp
trình bày bổ sung. & Phù hợp cho cả người đọc lẫn phần mềm xử lý
tiếp. \\
Mốc đến cuối tháng 6 & Ưu tiên parser ELF, resolver cơ bản, decoder
RV32I và JSON schema ổn định. & Đây là mốc khả thi, có thể demo rõ ràng
và tạo nền cho các chương sau. \\
\end{longtable}

\textbf{CHƯƠNG 3. PHÂN TÍCH THIẾT KẾ HỆ THỐNG}

Nội dung chương này sẽ được bổ sung ở giai đoạn tiếp theo. Trọng tâm dự
kiến gồm phân tích thiết kế phần mềm, giải thích cấu trúc từng lớp, mô
tả package, class chính và sơ đồ phụ thuộc giữa các thành phần.

\textbf{CHƯƠNG 4. HƯỚNG DẪN TRIỂN KHAI VÀ CÀI ĐẶT}

Nội dung chương này sẽ được bổ sung ở giai đoạn tiếp theo. Trọng tâm dự
kiến gồm yêu cầu môi trường, quy trình build, cách chạy chương trình,
cấu hình tham số đầu vào và các bước kiểm thử cài đặt cơ bản.

\textbf{CHƯƠNG 5. KẾT LUẬN}

Nội dung chương này sẽ được bổ sung ở giai đoạn tiếp theo. Phần hoàn
thiện sẽ tổng kết kết quả đạt được, hạn chế hiện tại và định hướng mở
rộng của hệ thống.

\textbf{TÀI LIỆU THAM KHẢO}

{[}1{]} RISC-V International, "The RISC-V Instruction Set Manual, Volume
I: Unprivileged Architecture," docs.riscv.org. {[}Online{]}. Available:
https://docs.riscv.org/reference/isa/\_attachments/riscv-unprivileged.pdf.
{[}Accessed: Mar. 29, 2026{]}.

{[}2{]} Free Software Foundation, "GNU Binutils Manual," sourceware.org.
{[}Online{]}. Available: https://sourceware.org/binutils/docs/binutils/.
{[}Accessed: Mar. 29, 2026{]}.

{[}3{]} Free Software Foundation, "BFD front end," sourceware.org.
{[}Online{]}. Available: https://sourceware.org/binutils/docs/bfd.pdf.
{[}Accessed: Mar. 29, 2026{]}.

{[}4{]} D. Andriesse, Practical Binary Analysis: Build Your Own Linux
Tools for Binary Instrumentation, Analysis, and Disassembly. San
Francisco, CA, USA: No Starch Press, 2019.

{[}5{]} Oracle, "Foreign Function and Memory API," Java SE 21
Documentation. {[}Online{]}. Available:
https://docs.oracle.com/en/java/javase/21/core/foreign-function-and-memory-api.html.
{[}Accessed: Mar. 29, 2026{]}.

{[}6{]} Apache Software Foundation, "What is Maven?," maven.apache.org.
{[}Online{]}. Available: https://maven.apache.org/what-is-maven.html.
{[}Accessed: Mar. 29, 2026{]}.

{[}7{]} Linux Foundation, "Chapter 4: Object Files," System V ABI.
{[}Online{]}. Available:
https://refspecs.linuxfoundation.org/elf/gabi4+/ch4.intro.html.
{[}Accessed: Mar. 29, 2026{]}.

{[}8{]} Linux Foundation, "Relocation" and "Symbol Table," System V ABI.
{[}Online{]}. Available:
https://refspecs.linuxfoundation.org/elf/gabi4+/ch4.reloc.html and
https://refspecs.linuxfoundation.org/elf/gabi4+/ch4.symtab.html.
{[}Accessed: Mar. 29, 2026{]}.

{[}9{]} National Security Agency, "P-Code Reference Manual," ghidra.re.
{[}Online{]}. Available:
https://ghidra.re/ghidra\_docs/languages/html/pcoderef.html.
{[}Accessed: Mar. 29, 2026{]}.

{[}10{]} Capstone Engine, "Documentation," capstone-engine.org.
{[}Online{]}. Available:
https://www.capstone-engine.org/documentation.html. {[}Accessed: Mar.
29, 2026{]}.

{[}11{]} N. Jones, "How to Create Jump Tables via Function Pointer
Arrays in C and C++," \emph{Netrino}, May 1999. {[}Online{]}. Available:
. {[}Accessed: May. 07, 2026{]}.

\end{document}
